\documentclass[professionalfont, 12pt, dvipdfmx, default, cjk]{beamer}
% \usetheme{Antibes}
% \usecolortheme{beaver}

\usepackage{tikz}
\usetikzlibrary{calc}
\usepackage[none]{hyphenat}
\usepackage{longtable}
\usepackage{booktabs}
\usepackage{graphicx}
\usepackage{amsmath, amssymb, cite, url}
\usepackage{etex}
\usepackage{longtable}

\setbeamercovered{dynamic} % slightly displaying not highlight

% font
\renewcommand{\kanjifamilydefault}{\gtdefault}
\mathversion{bold}
\usepackage{mathpazo}

\makeatletter
\def\Hline{%
  \noalign{\ifnum0=`}\fi\hrule \@height 0.8pt \futurelet
\reserved@a\@xhline}
\makeatother

\def\emph#1{\alert{#1}}

% define colours
% taken from pickton on Adobe Kuler:
% https://kuler.adobe.com/Some-Kind-Of-Ponde-color-theme-3837185/
\definecolor{PondeRed}{RGB}{230,37,52}
\definecolor{PondeWine}{RGB}{230,3,5}
\definecolor{PondePink}{RGB}{230,64,104}
\definecolor{PondeBlack}{RGB}{43,40,40}
\definecolor{PondeBlue}{RGB}{22,190,207}
\definecolor{PondeWhite}{RGB}{255,255,243}
\definecolor{PondeGray}{RGB}{107,110,108}

% set colours
\setbeamercolor{structure}{fg=PondeRed}
\setbeamercolor{itemize item}{fg=PondeRed}
\setbeamercolor{enumerate item}{fg=PondeRed}
\setbeamercolor{alerted text}{fg=PondeWine}
\setbeamercolor{section in toc}{fg=PondeBlack}

% set fonts
\setbeamerfont{itemize/enumerate body}{size=\normalsize}
\setbeamerfont{itemize/enumerate subbody}{size=\normalsize}
\setbeamerfont{itemize/enumerate subsubbody}{size=\small}

% item bullet
\setbeamertemplate{itemize item}{
  \tikz{
    \draw[fill=PondeWine,draw=none] (0, 0) rectangle(0.1, 0.1);
    \draw[fill=PondeWine,draw=none] (0.1, 0.1) rectangle(0.2, 0.2);
    \draw[fill=PondeWine,draw=none] (0, 0.2) rectangle(0.1, 0.3);
  }
}
\setbeamertemplate{itemize subitem}{
  \tikz{
    \draw[fill=PondeRed,draw=none] (0, 0) rectangle(0.075, 0.075);
    \draw[fill=PondeRed,draw=none] (0.075, 0.075) rectangle(0.15, 0.15);
    \draw[fill=PondeRed,draw=none] (0, 0.15) rectangle(0.075, 0.225);
  }
}
\setbeamertemplate{itemize subsubitem}{
  \tikz{
    \draw[fill=PondePink,draw=none] (0, 0) rectangle(0.075, 0.075);
    \draw[fill=PondePink,draw=none] (0.075, 0.075) rectangle(0.15, 0.15);
    \draw[fill=PondePink,draw=none] (0, 0.15) rectangle(0.075, 0.225);
  }
}

% title and background
\setbeamertemplate{title page}{}
\setbeamertemplate{background}{
  \begin{tikzpicture}
    \useasboundingbox (0,0) rectangle (\the\paperwidth,\the\paperheight);

    % background (solid=White)
    \fill[color=PondeWhite] (0,0) rectangle(\the\paperwidth,\the\paperheight);

    % draw title page as background
    \ifnum\thepage=1\relax
      \fill[color=PondeRed] (0, \paperheight-8mm) rectangle(\paperwidth,\paperheight);
      \fill[color=PondeBlack] (0, \paperheight-4mm) rectangle(\paperwidth,\paperheight);
      \node[anchor=south,text width=\paperwidth-1cm,inner xsep=0.5cm] at (0.5\paperwidth,4cm) {\color{PondeRed}\Huge\textbf{\inserttitle}};
      \node[anchor=west,inner xsep=0.5cm] at (0,2cm) {\color{PondeBlack}\LARGE\insertauthor};
      \node[anchor=south east] at(\paperwidth,0cm) {\color{PondeGray}\tiny\insertdate};
    \fi
  \end{tikzpicture}
}

% header
\makeatletter
\setbeamertemplate{headline}{}
\setbeamercolor{frametitle}{fg=PondeWhite,bg=PondeRed}
\setbeamertemplate{frametitle}
{
  \vspace*{-1mm}
  \begin{beamercolorbox}[sep=2mm,wd=1.01\paperwidth]{frametitle}
    \hfill \tiny \insertsectionhead \par
    \vspace*{-2mm} \Large \strut\insertframetitle\strut
  \end{beamercolorbox}
}

\makeatother

% footer is progress bar
\makeatother
\setbeamertemplate{footline}
{
  \dimen0=1.1\paperwidth
  \newcount\p
  \newcount\q
  \p=\insertframenumber
  \q=\inserttotalframenumber
  \advance\p by -1
  \advance\q by -1
  \ifnum\q=0
    \q=1
  \fi
  \multiply\dimen0 by \the\p
  \divide\dimen0 by \the \q
  \edef\progressbarwidth{\the\dimen0}
  \leavevmode%
  \setbeamercolor{progress outer}{fg=white, bg=PondeGray}
  \setbeamercolor{progress bar}{bg=PondeRed}
  \begin{beamercolorbox}[wd=1.1\paperwidth,ht=0.25ex,dp=1ex]{progress outer}
    \begin{beamercolorbox}[wd=\progressbarwidth,ht=0.25ex,dp=1ex]{progress bar}
    \end{beamercolorbox}%
  \end{beamercolorbox}%
}
\makeatletter

% page number
\setbeamertemplate{navigation symbols}{
  \setbeamercolor{footline}{fg=PondeGray}
  \usebeamerfont{footline}
  \usebeamercolor[fg]{footline}
  \ifnum\thepage>1\relax
    \tiny \insertframenumber / \inserttotalframenumber
    \hspace{0.2em}
  \fi
}

% Comment these out if you don't want a slide with just the
% part/section/subsection/subsubsection title:
% \AtBeginPart{\frame{\partpage}}
% \AtBeginSection{\frame{\sectionpage}}
% \AtBeginSubsection{\frame{\subsectionpage}}
% \AtBeginSubsubsection{\frame{\subsubsectionpage}}
\setlength{\parindent}{0pt}
% \setlength{\parskip}{6pt plus 2pt minus 1pt}
\setlength{\emergencystretch}{3em}  % prevent overfull lines
\setcounter{secnumdepth}{0}

\title{言語の極限同定みたいな話}
\author{Hattori Kazuhiro}
\date{2015/11/13}

\InputIfFileExists{macro.tex}{}{}

\begin{document}
\frame{\titlepage}

\section{Reference}\label{reference}

\begin{frame}{参考文献}

\begin{itemize}
\itemsep1pt\parskip0pt\parsep0pt
\item
  \href{https://en.wikipedia.org/wiki/Language_identification_in_the_limit}{wikipedia/Language\_identification\_in\_the\_limit}
\item
  \href{https://en.wikipedia.org/wiki/Grammar_induction}{wikipedia/Grammar\_induction}
\item
  \href{http://web.mit.edu/6.863/www/spring2010/readings/gold67limit.pdf}{paper/gold67limit.pdf}
\item
  \href{http://www.sciencedirect.com/science/article/pii/0022000080900410}{paper/Angluin80}
\item
  \href{http://repository.kulib.kyoto-u.ac.jp/dspace/bitstream/2433/103408/1/0482-7.pdf}{paper/Shinohara83.pdf}
\item
  \href{http://www-ikn.ist.hokudai.ac.jp/~arim/papers/arimura_stacs94.pdf}{paper/Arimura94.pdf}
\end{itemize}

\end{frame}

\section{Definitions}\label{definitions}

\begin{frame}{諸定義}

\begin{itemize}
\itemsep1pt\parskip0pt\parsep0pt
\item
  アルファベットとはシンボルの\alert{有限の}集合

  \begin{itemize}
  \itemsep1pt\parskip0pt\parsep0pt
  \item
    文脈によってはシンボルは2つ以上に限ることがある
  \item
    \(\Sigma = \{ 0, 1, \cdots \}\)
  \end{itemize}
\item
  テキストとはアルファベットの上の文字列

  \begin{itemize}
  \itemsep1pt\parskip0pt\parsep0pt
  \item
    \(\Sigma^* = \{ \epsilon, 0, 1, \cdot, 00, 01, 10, 11, \cdots, 010, \cdots \}\)

    \begin{itemize}
    \itemsep1pt\parskip0pt\parsep0pt
    \item
      空文字列を \(\epsilon\) と書く
    \end{itemize}
  \end{itemize}
\item
  言語とはテキストの集合

  \begin{itemize}
  \itemsep1pt\parskip0pt\parsep0pt
  \item
    \(L \subseteq \Sigma^*\)
  \end{itemize}
\end{itemize}

\end{frame}

\section{E. Mark Gold}\label{e.-mark-gold}

\begin{frame}{E. Mark Gold:}

\fontsize{19pt}{20pt} \selectfont

``Language Identification in the Limit'', 1967

\end{frame}

\begin{frame}{動機}

\begin{itemize}
\itemsep1pt\parskip0pt\parsep0pt
\item
  AIに言葉を話させたい (作文をさせたい)
\item
  言葉を話すためにはその言語を学習しなければならない

  \begin{enumerate}
  \def\labelenumi{\arabic{enumi}.}
  \itemsep1pt\parskip0pt\parsep0pt
  \item
    AI (計算機) は言語を学習できるか
  \item
    どの時点で学習したと言えるか
  \end{enumerate}
\end{itemize}

\end{frame}

\begin{frame}{言葉を話せるAI}

\begin{itemize}
\itemsep1pt\parskip0pt\parsep0pt
\item
  人間の場合は、英語のルールを正しく書き下せなくても、``英語を話せる''
  と言う
\item
  AI の場合でも、同様の ``英語を話せる''
  とかいえるようなモデルが作れるはず (模倣)

  \begin{itemize}
  \item
    すなわち言語のルールを書き下すことなく学習ができるはず \pause
  \item
    英語に限らず自然言語一般のルールを学習させたい

    \begin{itemize}
    \itemsep1pt\parskip0pt\parsep0pt
    \item
      人間は任意の自然言語を学習可能
    \end{itemize}
  \item
    自然言語をよりも単純な言語はより用意に学習できる

    \begin{itemize}
    \itemsep1pt\parskip0pt\parsep0pt
    \item
      自然言語を lower bound とする
    \item
      より複雑な言語は別にどうでもいい
    \end{itemize}
  \item
    少なくとも自然言語を含むような単純な言語の枠組みを設定する
  \end{itemize}
\end{itemize}

\end{frame}

\begin{frame}{言語の枠組み: 言語クラス}

何の枠組みもなしに言語を学習するってのは難しすぎる

\vspace*{3mm}

\centering
言語クラス \(C = \{ L_1, L_2, \cdots \}\)
\footnote{個別の言語について、それが学習可能か？ を議論せず、言語クラスに対して議論する}\\\footnotesize \textcolor{gray}{(実際は添字付きとは限らない)}
\normalsize
\flushleft

\begin{itemize}
\itemsep1pt\parskip0pt\parsep0pt
\item
  \(C\) はありえる言語全体
  \footnote{学習したい言語が正規言語だと分かってるとき、正規言語全体というクラスを指定すればいい}

  \begin{itemize}
  \itemsep1pt\parskip0pt\parsep0pt
  \item
    例えば「ありえる自然言語全体」を設定する
  \item
    例えば「文脈自由文法全体」を設定する

    \begin{itemize}
    \itemsep1pt\parskip0pt\parsep0pt
    \item
      その中には実在する自然言語 (英語など) が含まれる
    \end{itemize}
  \end{itemize}
\item
  言語の学習とは言語クラス \(C\) からそれっぽい言語を選ぶことだとする
\end{itemize}

\end{frame}

\begin{frame}{情報提示 (information presentation)}

何はともあれ、情報を提示されないと学習できない

Gold は Text 、Informant の2つがあると考えた

現代語訳すると

\begin{itemize}
\itemsep1pt\parskip0pt\parsep0pt
\item
  Text \(\Rightarrow\) 正提示
\item
  Informant \(\Rightarrow\) 完全提示
\end{itemize}

紛らわしいので現代語を使います

\end{frame}

\begin{frame}{情報提示/正提示}

\begin{itemize}
\itemsep1pt\parskip0pt\parsep0pt
\item
  正提示: 言語に出現するテキストの (無限) 列 \vspace*{-2mm}
  \[w_1, w_2, \cdots\]

  \begin{itemize}
  \itemsep1pt\parskip0pt\parsep0pt
  \item
    言語に含まれるテキストはこの列のどこかで必ず出現することが保証される:
    (\(\forall w \in L, \exists i \in \mathbb{N}, w_i = w\)) \footnote{
    謎の操作 $content$ を使って次のようにも書く:
    $$content(w_1, w_2, \cdots) = L$$
    }
  \end{itemize}
\end{itemize}

\pause
\small

\begin{itemize}
\itemsep1pt\parskip0pt\parsep0pt
\item
  これは次に相当する

  \begin{itemize}
  \itemsep1pt\parskip0pt\parsep0pt
  \item
    子供が、大人から常に正しいテキストを一つずつ聞く
  \end{itemize}
\end{itemize}

\end{frame}

\begin{frame}{情報提示/完全提示}

\begin{itemize}
\item
  完全提示:
  何でもいいから任意のテキストの列で、そのテキストが言語に含まれるかのラベル
  (\texttt{Bool}) が付いている \[(x_1, I_1), (x_2, I_2), \cdots\]
\item
  \(I_i \land (x_i \in L)\) \hspace{1mm} or \hspace{1mm}
  \(\lnot I_i \land (x_i \not\in L)\)
\item
  ただし \(content(x_1, x_2, \cdots) = \Sigma^*\)
\end{itemize}

\pause
\small

\begin{itemize}
\itemsep1pt\parskip0pt\parsep0pt
\item
  次に相当

  \begin{itemize}
  \itemsep1pt\parskip0pt\parsep0pt
  \item
    子供が大人から常に正しいテキストを聞く (\(I_i = true\))
  \item
    子供はたまに作文をする (話す)

    \begin{itemize}
    \itemsep1pt\parskip0pt\parsep0pt
    \item
      文法エラーをすると、リアクションから誤りだと知る (\(I_i = false\))
    \end{itemize}
  \end{itemize}
\end{itemize}

\end{frame}

\begin{frame}{言語学習モデル (guessing machine)}

\begin{enumerate}
\def\labelenumi{\arabic{enumi}.}
\itemsep1pt\parskip0pt\parsep0pt
\item
  どちらかの提示により、学習者は \(i_1, i_2, \cdots\) を逐次的に受け取る
\item
  学習者は謎の手続き \(G\) によって言語を推論する
\end{enumerate}

時刻 \(t\) における推論: \[g_t = G(i_1, i_2, \cdots, i_t)\]

\(g_t\) の型は特定の一つの言語を説明するもの、或いは、言語そのもの
(\(L(g_t) \in C\) 或いは \(g_t \in C\))

\end{frame}

\begin{frame}{極限における同定}

原理的に有限個の情報からは学習できない、という言語クラスはいくらでもあるので極限を考える

\[g_t = G(i_1, i_2, \cdots, i_t)\]

が極限 \(t \rightarrow \infty\) で同定してかつ正しい言語を指し示すことを
「極限同定」 という

\begin{itemize}
\itemsep1pt\parskip0pt\parsep0pt
\item
  \emph{N.B.}
  学習者は「いつ自分は正しく言語を学習できたか」を知る必要はないし、知る術はない
\end{itemize}

\end{frame}

\begin{frame}{言語クラスの学習可能性}

Gold が提示した枠組みにおける「学習可能性」とは次のこと

\vspace*{4mm}

\begin{itemize}
\itemsep1pt\parskip0pt\parsep0pt
\item
  「ある言語クラスが学習可能」を示すためには、

  \begin{itemize}
  \itemsep1pt\parskip0pt\parsep0pt
  \item
    ある guessing machine を構成して、
  \item
    その推論結果が常に極限で同定できていることを言えば良い
  \end{itemize}
\end{itemize}

\end{frame}

\begin{frame}{完全提示から正規言語は学習可能}

\begin{itemize}
\itemsep1pt\parskip0pt\parsep0pt
\item
  言語からオートマトン (\(g_t\)) を復元する
\item
  言語から正規表現を得られる

  \begin{itemize}
  \itemsep1pt\parskip0pt\parsep0pt
  \item
    但し正例と負例を無限に作る必要がある

    \begin{itemize}
    \itemsep1pt\parskip0pt\parsep0pt
    \item
      例えば、任意のテキストについて、言語に含まれるかどうか判断する神託
    \item
      あるいは手元にあるオートマトンと同値なオートマトンを作りたい?
    \end{itemize}
  \end{itemize}
\end{itemize}

\end{frame}

\begin{frame}{libalf: Automata Learning Framework}

オートマトンを復元してくれるライブラリ:

libalf
(\href{http://libalf.informatik.rwth-aachen.de/index.php?page=about}{libalf.informatik.rwth-aachen.de})
とかゆう実装がある

\vspace*{3mm}

\centering
\includegraphics[width=1.00\textwidth,bb=0 0 1294 412]{img/libalf.jpg}

\end{frame}

\begin{frame}{libalf/demo/online}

\includegraphics[width=1.0\textwidth,bb=0 0 1667 1297]{img/online1.jpg}

\end{frame}

\begin{frame}{out}

\includegraphics[width=1.0\textwidth,bb=0 0 304 112]{img/1.png}

\vfill
\pause
僕「\(0^n1^n\) を受理させたい!!」\\上は\(11\) を受理している! おかしい!!

\end{frame}

\begin{frame}{cont}

\centering
\includegraphics[width=1.0\textwidth,bb=0 0 1472 607]{img/online2.jpg}

\end{frame}

\begin{frame}{out}

\includegraphics[width=1.0\textwidth,bb=0 0 413 151]{img/2.png}

\end{frame}

\begin{frame}{out}

\includegraphics[width=1.0\textwidth,bb=0 0 752 241]{img/3.png}

\end{frame}

\begin{frame}{out}

\includegraphics[width=1.0\textwidth,bb=0 0 768 231]{img/4.png}

\vfill
\(0^n1^n\) は正規言語じゃない?!?!?

\end{frame}

\begin{frame}{\texttt{/(000)*/}}

\includegraphics[width=1.0\textwidth,bb=0 0 413 116]{img/5.png}

\vfill
\(\epsilon, 0, 00, 000, 0000, 00000\) を神託した時点で上を出力

\end{frame}

\begin{frame}{正提示から正規言語を極限同定することは不可能}

\begin{itemize}
\itemsep1pt\parskip0pt\parsep0pt
\item
  正規言語全体は大きすぎる (超有限) ため
\item
  どの時刻 \(t \in \mathbb{N}\) においても一つの言語に絞ることができない
\end{itemize}

\vfill
Gold 曰く、``正提示は完全提示に較べて弱い''

\end{frame}

\begin{frame}{子供の言語学習 (acquisition of grammar by children)}

ちょいちょい「子供が学習するときは〜」という話が出るが、実際はどうなのか

\vspace*{3mm}

psycholinguist 曰く {[}McNeill,
1966{]}:\\「子供が文法誤りをしたとき、それを指摘することは滅多にない」

\flushleft
\vspace*{3mm}

\begin{itemize}
\itemsep1pt\parskip0pt\parsep0pt
\item
  完全提示ってのはちょっと仮定が強すぎるのでは？

  \begin{itemize}
  \itemsep1pt\parskip0pt\parsep0pt
  \item
    でも我々は自然言語を学習する
  \end{itemize}
\end{itemize}

\end{frame}

\begin{frame}{自然言語は正提示から学習可能説}

\begin{itemize}
\itemsep1pt\parskip0pt\parsep0pt
\item
  多くの自明な言語クラスなら正提示から学習可能であることを Gold は示した
\item
  英語は文脈自由文法だと言われているが、実際には、全ての文脈自由文法が自然言語になるわけではない

  \begin{itemize}
  \itemsep1pt\parskip0pt\parsep0pt
  \item
    もっと制限がある (もっと小さなクラスである)
  \item
    例えば、学習可能性の結論として:
    あり得る自然言語のクラスは、無限言語を少なくともひとつは含み、
    全ての有限言語を含むことはない (超有限ではない)
  \end{itemize}
\end{itemize}

\end{frame}

\begin{frame}{子供は、我々がわからない方法で負例を学習する説}

\begin{itemize}
\itemsep1pt\parskip0pt\parsep0pt
\item
  例えば、発言をして、思ったような反応が得られなかったとき
\item
  だとすると完全提示からの学習をしてもいい
\item
  原始再帰的言語は完全提示から学習可能であることをGoldは示している

  \begin{itemize}
  \itemsep1pt\parskip0pt\parsep0pt
  \item
    文脈依存文法も原始再帰的言語の一つ
  \item
    英語は完全提示から学習可能
  \end{itemize}
\end{itemize}

\end{frame}

\section{Dana Angluin}\label{dana-angluin}

\begin{frame}{Dana Angluin}

\fontsize{19pt}{20pt} \selectfont
``On the Complexity of Minimum Inference of Regular Sets'', 1978

``Finding Patterns Common to a Set of Strings'', 1979

\end{frame}

\begin{frame}{Angluin さんのやったこと}

正提示から学習可能な非自明な言語クラスの発掘

\begin{enumerate}
\def\labelenumi{\arabic{enumi}.}
\itemsep1pt\parskip0pt\parsep0pt
\item
  パターン言語
\end{enumerate}

\end{frame}

\begin{frame}{パターン: \(\mathcal{PAT}\)}

\begin{itemize}
\item
  有限アルファベット \(\Sigma = \{ 0, 1, \cdots \}\)
\item
  無限変数 \(X = \{ x_1, x_2, \cdots \}\)
  \[\mathcal{PAT} = (\Sigma \cup X)^+\]
\item
  代入、\(\preceq\) (less-general-than)

  \begin{itemize}
  \itemsep1pt\parskip0pt\parsep0pt
  \item
    代入 \([x_i/p]\) は、パターン中に出現する\alert{全ての} 変数 \(x_i\)
    を \alert{空でない} パターン \(p\) に置き換える操作
  \end{itemize}
\end{itemize}

\vspace*{-6mm}\begin{align*}
x_1 x_1 & \succeq \underline{x_2} ~ \underline{x_2} & [x_1/x_2] \\
& \succeq \underline{a x_1} ~ \underline{a x_1} & [x_2/a x_1] \\
& \succeq a \underline{b c} a \underline{b c} & [x_1/b c] \\
\end{align*}

\vspace*{-10mm}

\[a x_1 a x_1 \not\succeq a a\]

\end{frame}

\begin{frame}{パターン言語}

\[L(p) = \{ w \in \Sigma^* | w \preceq p \}\]

\vfill
例えば、 \[L(x_1 x_1) = \{ w w | w \in \Sigma^+ \}\]

\end{frame}

\begin{frame}{パターン言語は正提示から学習可能}

パターン言語のクラス \[C = \{ L(p) | p \in \mathcal{PAT} \}\]

\end{frame}

\begin{frame}{学習過程 (推論機械 \(G\))}

\begin{itemize}
\itemsep1pt\parskip0pt\parsep0pt
\item
  \(L = L(p) \in C\) を学習したい
\item
  正提示 \(s_1, s_2, \cdots (s_i \in L)\) を受け取る
\end{itemize}

推論機械はパターンを出力する: \[p_t = G(s_1 \cdots s_t)\]

\pause
大雑把にいうと

\begin{itemize}
\itemsep1pt\parskip0pt\parsep0pt
\item
  \(\forall t, content(s_1, \cdots, s_t) \subseteq L(p_t)\) (無矛盾)
\item
  \(L(p_1) \subseteq L(p_2) \subseteq \cdots\) (保守的)
\item
  上2つを満たす為には次のようにすれば十分 (極小言語戦略)

  \begin{itemize}
  \itemsep1pt\parskip0pt\parsep0pt
  \item
    \(p_t = \mathop{\rm arg\;min}\limits_p L(p) ~s.t.~ content \subseteq L(p)\)
  \end{itemize}
\end{itemize}

ってやると、 \[L = L(\lim_t p_t)\] になる.

\end{frame}

\begin{frame}{皆が興味がなさそうだったら飛ばすページ}

\begin{block}{Prop.
有限の厚みを持つ言語クラスは極小言語によって極限同定可能である}

言語クラスが有限の厚みを持つとは \(\iff\) 任意のテキストの有限集合 \(S\)
について \(\{ L | S \subseteq L \}\) が有限であること

\begin{enumerate}
\def\labelenumi{\arabic{enumi}.}
\itemsep1pt\parskip0pt\parsep0pt
\item
  無矛盾かつ保守的な推論による推論の列: \(g_1, g_2, ...\) は収束する

  \begin{enumerate}
  \def\labelenumii{\arabic{enumii}.}
  \itemsep1pt\parskip0pt\parsep0pt
  \item
    有限の厚みより、正提示の1つ目を含む言語は有限個しかない

    \begin{enumerate}
    \def\labelenumiii{\arabic{enumiii}.}
    \itemsep1pt\parskip0pt\parsep0pt
    \item
      推論も有限個しかない
    \end{enumerate}
  \item
    保守性より \(L(g_1) \subseteq L(g_2) \subseteq \cdots\) (\(g_i\)
    に半順序がつく)
  \item
    推論列はどこかで停まるか、極大を定める
  \end{enumerate}
\end{enumerate}

\end{block}

\end{frame}

\begin{frame}{パターン言語は正提示から学習可能であることの証明}

\begin{enumerate}
\def\labelenumi{\arabic{enumi}.}
\itemsep1pt\parskip0pt\parsep0pt
\item
  先に挙げた極小言語戦略による推論機械を構成する
\item
  パターン言語クラスは有限の厚みを持つことを示す

  \begin{itemize}
  \itemsep1pt\parskip0pt\parsep0pt
  \item
    \(\forall w, \{ p | w \preceq p \}\) が有限であることを言えばよい
  \item
    hint: \(q \preceq p \Rightarrow |q| \geq |p|\)
  \end{itemize}
\item
  先の Prop. から極限同定可能
\end{enumerate}

\end{frame}

\section{Others}\label{others}

\begin{frame}{Shinohara: ``Polynomial Time Inference of Extended Regular
Pattern Languages'', 1991}

パターン言語を消去可能パターン言語に拡張

\begin{itemize}
\itemsep1pt\parskip0pt\parsep0pt
\item
  空の代入を許す: \[a x_1 a x_1 \succeq a a\]
\end{itemize}

消去可能パターン言語であって正則なものは、正提示から学習可能であることを示した

\begin{itemize}
\itemsep1pt\parskip0pt\parsep0pt
\item
  N.B. パターン言語が正則 \(\iff\) 一つのパターンに出現する同じ変数
  \(x_i\) は高々一つ (出現する変数が全て異なる)
\end{itemize}

\end{frame}

\begin{frame}{使い道}

\includegraphics[width=1.0\textwidth,bb=0 0 1186 794]{img/shinohara.jpg}

\end{frame}

\begin{frame}{Arimura, Shinohara: ``Finding Minimal Generalizations of
Unions of Pattern Languages and \ldots{}'', 1994}

パターンの和言語も正提示から学習可能であることを示した

\begin{itemize}
\itemsep1pt\parskip0pt\parsep0pt
\item
  大体一つのパターンの言語なんて単純すぎる
\item
  \(L(p_1, \cdots, p_k) = L(p_1) \cup \cdots \cup L(p_k)\)

  \begin{itemize}
  \itemsep1pt\parskip0pt\parsep0pt
  \item
    ただしパターンの数に上限 \(k\) を設定する
  \end{itemize}
\end{itemize}

\end{frame}

\begin{frame}{Takeuchi, Sato:
``誤情報を含む正則パターン言語の多項式時間推論'', 1998}

タイトルの通り

\begin{itemize}
\itemsep1pt\parskip0pt\parsep0pt
\item
  学習したい言語からちょっとくらいズレた提示があっても良いようにする
\item
  言語の「近傍系」を学習していく

  \begin{itemize}
  \itemsep1pt\parskip0pt\parsep0pt
  \item
    パターンの上の一種のハミング距離を取る
  \end{itemize}
\end{itemize}

(一定の) 誤情報を含む正提示から学習可能

\end{frame}

\begin{frame}{Ng, Shinohara: ``Inferring Unions of the Pattern Languages
by the Most Fitting Covers'', 2005}

\begin{itemize}
\itemsep1pt\parskip0pt\parsep0pt
\item
  今見てきた極小言語戦略は提示 \(s_1, \cdots, s_t\)
  を含む中で極小な言語を推論として出力する

  \begin{itemize}
  \itemsep1pt\parskip0pt\parsep0pt
  \item
    極小とは言語の包含関係に関する
  \item
    言語は無限集合なので大きさを比較するのは難しい
  \end{itemize}
\item
  要するに小さければ良い

  \begin{itemize}
  \itemsep1pt\parskip0pt\parsep0pt
  \item
    実際の言語ではテキストの長さには上限が (たぶん) ある

    \begin{itemize}
    \itemsep1pt\parskip0pt\parsep0pt
    \item
      \(\Rightarrow\) 任意の自然言語は有限集合である
    \item
      大きさに関する極小を取ることができる
    \end{itemize}
  \end{itemize}
\end{itemize}

性質を調べたものであって、この場合の推論方法を示したわけでも学習可能性を言ったものでもない

\end{frame}

\section{おわりに}\label{ux304aux308fux308aux306b}

\begin{frame}{おわりに}

\begin{itemize}
\itemsep1pt\parskip0pt\parsep0pt
\item
  誰も自然言語のことを忘れてパターン言語だけをやっている
\item
  完全提示からの学習なんて誰もやってない (ことはないけど)
\item
  というか確率を導入すべきだ
\item
  もはや誰も手をつけない分野

  \begin{itemize}
  \itemsep1pt\parskip0pt\parsep0pt
  \item
    ちょっと古ければ論文は大量にある
  \item
    2000年以降もたま〜にぽつぽつ出てる
  \end{itemize}
\end{itemize}

\end{frame}

\end{document}
